\chapter{Software Requirements Specification}

\section{Introduction}

\subsection{Document Purpose}
This document specifies the software requirements for the \textbf{LivingDocs -- Design and Implementation of an AI-Assisted Self-Updating Document System for Engineering Teams} platform.

LivingDocs is an AI-powered documentation maintenance platform engineered to detect documentation drift, automatically generate and update technical documentation from source code, and enforce a multi-tiered Human-in-the-Loop review workflow.

This specification focuses on:
\begin{itemize}
	\item User requirements and target personas.
	\item Functional Requirements (FR).
	\item Non-Functional Requirements (NFR).
	\item User roles and access control mechanisms.
	\item Documentation drift detection and management workflows.
	\item Verification, auditing, and document versioning controls.
\end{itemize}

\textit{Note: This document excludes low-level architectural blueprints, database schema design, or specific code implementation details, which are addressed in subsequent system design modules.}

\section{System Scope}
LivingDocs automates the continuous maintenance of technical documentation across software engineering projects.

Key capabilities include:
\begin{itemize}
	\item Seamless integration with GitHub repositories via OAuth and Webhooks.
	\item Continuous ingestion of commit updates and Pull Requests.
	\item Automated static code and AST analysis.
	\item Bidirectional mapping between code entities and documentation artifacts.
	\item Multi-classification of documentation drift (Referential, Signature, Semantic).
	\item AI-assisted document generation and context-aware auto-updates.
	\item Granular version history and complete audit traceability.
	\item Multi-tiered review workflows involving Staff Reviewers and Manager Approvers.
	\item Searchable Knowledge Base linking documentation directly to live code.
	\item System-wide governance and administrative controls.
\end{itemize}

\section{System Objectives}
The primary objectives of LivingDocs encompass:
\begin{itemize}
	\item Minimizing technical documentation debt in fast-paced development environments.
	\item Proactively identifying outdated or contradictory documentation.
	\item Detecting structural reference breaks as well as semantic behavioral shifts.
	\item Reducing manual overhead required for writing and updating technical docs.
	\item Ensuring high precision and global consistency across project documentation.
	\item Maintaining strict human governance over AI-generated outputs.
	\item Guaranteeing complete auditability and version rollback capabilities.
\end{itemize}

\section{Target User Personas}
The system defines five primary operational roles:

\begin{longtable}{p{4.5cm} p{9cm}}
	\toprule
	\textbf{Role} & \textbf{Description} \\
	\midrule
	\endhead
	
	Developer (Author) &
	Engineers who connect repositories, initiate AI documentation generation, review drift alerts, and submit documentation drafts for review. \\
	
	Staff (Reviewer) &
	Technical personnel responsible for conducting the initial technical review of AI-generated or auto-updated documentation. \\
	
	Manager (Approver) &
	Decision-makers responsible for final approval, publication, or direct repository commit authorization. \\
	
	Technical Lead / Owner &
	Administrators managing documentation standards, drift policies, custom templates, and quality compliance. \\
	
	Administrator &
	System administrators managing user accounts, RBAC policies, GitHub integration setups, and backend infrastructure. \\
	
	\bottomrule
\end{longtable}

\section{User Requirements}

\subsection{Developer (Author)}
Developers require the capability to:
\begin{itemize}
	\item Authenticate via credentials or GitHub OAuth.
	\item Manage personal profiles and notification preferences.
	\item Link GitHub accounts and select repositories for workspace ingestion.
	\item Monitor repository sync status and trigger manual re-synchronization.
	\item View ingested commits and Pull Requests.
	\item Generate technical documentation from source code via AI.
	\item Select customized documentation templates and preview draft outputs.
	\item Identify code entities lacking documentation coverage.
	\item Edit AI-generated drafts and inspect grounding evidence before submission.
	\item Submit drafts into review workflows and monitor progress.
	\item Inspect documentation drift alerts, severity levels, and associated diffs.
	\item Compare document versions and perform version rollbacks.
	\item Search the knowledge base for linked code entities and documentation.
\end{itemize}

\subsection{Staff (Reviewer)}
Staff members require the capability to:
\begin{itemize}
	\item Access a centralized review queue.
	\item Inspect documentation side-by-side with code grounding evidence.
	\item Compare proposed versions against the latest approved baseline.
	\item Verify compliance with designated documentation templates and code truth.
	\item Inline review comments, request modifications, or make minor corrections.
	\item Approve drafts for Manager escalation or reject drafts with actionable feedback.
	\item Filter, sort, and track historical review logs.
\end{itemize}

\subsection{Manager (Approver)}
Managers require the capability to:
\begin{itemize}
	\item Access high-priority approval queues containing Staff-reviewed docs.
	\item Review Staff feedback, change diffs, and complete audit trails.
	\item Grant final approval for automatic repository commits or public publishing.
	\item Reject submissions or mandate further revisions.
	\item Execute emergency rollbacks to prior approved versions.
	\item Configure approval policies based on drift severity (e.g., Critical drift).
	\item Delegate or reassign approval authority.
\end{itemize}

\subsection{Technical Lead / Documentation Owner}
Technical Leads require the capability to:
\begin{itemize}
	\item Organize and manage documentation collections and ownership.
	\item Define drift detection sensitivity thresholds and update policies.
	\item Monitor the Documentation Health Dashboard (coverage, freshness, drift metrics).
	\item View code-to-documentation dependency graphs.
	\item Export health compliance and audit reports.
	\item Create, clone, modify, and publish version-controlled documentation templates.
	\item Configure mandatory template headings, sections, and placeholders.
\end{itemize">}

\subsection{Administrator}
Administrators require the capability to:
\begin{itemize}
	\item Manage user accounts, roles, and granular RBAC permissions.
	\item Configure workspace configurations, GitHub Apps, OAuth credentials, and webhooks.
	\item Define supported programming languages and parser configurations.
	\item Set AI inference endpoints, token quotas, and vector indexing parameters.
	\item Monitor event processing pipelines, background jobs, and system health.
	\item Inspect system-wide administrative audit logs.
\end{itemize">}

\section{Functional Requirements}

\subsection{Account and Authentication}
\begin{longtable}{p{2cm} p{11.3cm}}
	\toprule
	\textbf{ID} & \textbf{Requirement Description} \\
	\midrule
	\endhead
	FR01 & The system shall allow users to register new accounts. \\
	FR02 & The system shall allow users to securely log in and log out. \\
	FR03 & The system shall support authentication via Email/Password and GitHub OAuth. \\
	FR04 & The system shall enable users to manage profiles and notification preferences. \\
	\bottomrule
\end{longtable}

\subsection{GitHub Integration}
\begin{longtable}{p{2cm} p{11.3cm}}
	\toprule
	\textbf{ID} & \textbf{Requirement Description} \\
	\midrule
	\endhead
	FR05 & The system shall allow developers to link GitHub accounts via OAuth2. \\
	FR06 & The system shall display requested access scopes and permissions. \\
	FR07 & The system shall allow users to select and import target repositories. \\
	FR08 & The system shall support disconnecting repositories and revoking access tokens. \\
	FR09 & The system shall display synchronization status, default branch, and last sync timestamp. \\
	FR10 & The system shall support manual trigger of repository re-synchronization. \\
	FR11 & The system shall track ingested commits and Pull Request events. \\
	FR12 & The system shall alert administrators upon Webhook or token expiration. \\
	\bottomrule
\end{longtable}

\subsection{AI Source-Code Reading \& Documentation Generation}
\begin{longtable}{p{2cm} p{11.3cm}}
	\toprule
	\textbf{ID} & \textbf{Requirement Description} \\
	\midrule
	\endhead
	FR13 & The system shall support doc generation at file, module, or repository levels. \\
	FR14 & The system shall analyze AST and code structure to synthesize documentation. \\
	FR15 & The system shall allow selecting predefined documentation templates. \\
	FR16 & The system shall identify code entities lacking documentation coverage. \\
	FR17 & The system shall provide a real-time preview of AI-generated drafts. \\
	FR18 & The system shall allow Developers to edit drafts prior to submission. \\
	FR19 & The system shall display grounding evidence (code references) for generated docs. \\
	FR20 & The system shall allow submitting generated drafts into the review workflow. \\
	FR21 & The system shall display the current workflow state of the document. \\
	\bottomrule
\end{longtable}

\subsection{Automatic Documentation Update}
\begin{longtable}{p{2cm} p{11.3cm}}
	\toprule
	\textbf{ID} & \textbf{Requirement Description} \\
	\midrule
	\endhead
	FR22 & The system shall detect affected documentation upon Pull Request or commit updates. \\
	FR23 & The system shall automatically generate updated document drafts and increment versions. \\
	FR24 & The system shall allow toggling auto-update features per repository or module. \\
	FR25 & The system shall provide side-by-side diff views between auto-updated drafts and published docs. \\
	FR26 & The system shall highlight specific triggering commits and diff chunks. \\
	FR27 & The system shall support accepting, editing, or dismissing auto-updated suggestions. \\
	FR28 & The system shall allow manual re-generation of documentation updates on demand. \\
	FR29 & The system shall support one-click rollbacks to pre-update states. \\
	\bottomrule
\end{longtable}

\subsection{Documentation Drift Detection and Review}
\begin{longtable}{p{2cm} p{11.3cm}}
	\toprule
	\textbf{ID} & \textbf{Requirement Description} \\
	\midrule
	\endhead
	FR30 & The system shall analyze Pull Request diffs to detect documentation drift. \\
	FR31 & The system shall present grounding evidence for every drift alert. \\
	FR32 & The system shall classify drift into Referential, Signature, or Semantic categories. \\
	FR33 & The system shall assign drift severity levels: Low, Medium, High, or Critical. \\
	FR34 & The system shall generate context-aware fix suggestions for detected drift. \\
	FR35 & The system shall support manual trigger of drift analysis for specific files or PRs. \\
	FR36 & The system shall maintain historical logs of drift occurrences and resolution actions. \\
	\bottomrule
\end{longtable}

\subsection{Change Logs and Version History}
\begin{longtable}{p{2cm} p{11.3cm}}
	\toprule
	\textbf{ID} & \textbf{Requirement Description} \\
	\midrule
	\endhead
	FR37 & The system shall record full change histories for every document. \\
	FR38 & The system shall allow comparing differences across any two arbitrary versions. \\
	FR39 & The system shall log acting user/role, timestamp, and trigger event for each version. \\
	FR40 & The system shall allow restoring or rolling back documents to previous versions. \\
	\bottomrule
\end{longtable}

\subsection{Knowledge Base}
\begin{longtable}{p{2cm} p{11.3cm}}
	\toprule
	\textbf{ID} & \textbf{Requirement Description} \\
	\midrule
	\endhead
	FR41 & The system shall provide full-text search across documentation and linked code entities. \\
	FR42 & The system shall render interactive dependency graphs between code and documentation. \\
	FR43 & The system shall allow bidirectional navigation between code files and docs. \\
	\bottomrule
\end{longtable}

\subsection{Review and Approval Workflows}
\begin{longtable}{p{2cm} p{11.3cm}}
	\toprule
	\textbf{ID} & \textbf{Requirement Description} \\
	\midrule
	\endhead
	FR44 & The system shall provide a dedicated Staff Review Queue. \\
	FR45 & Staff shall be able to inspect grounding evidence, change logs, and template rules. \\
	FR46 & Staff shall be able to approve, reject, inline-edit, or request revisions. \\
	FR47 & The system shall provide a Manager Approval Queue for Staff-approved documents. \\
	FR48 & Managers shall have final authority to publish or execute automated repository commits. \\
	FR49 & The system shall enforce mandatory approval rules before committing changes. \\
	FR50 & The system shall record complete, immutable audit logs of all approval decisions. \\
	\bottomrule
\end{longtable}

\subsection{Governance and Administration}
\begin{longtable}{p{2cm} p{11.3cm}}
	\toprule
	\textbf{ID} & \textbf{Requirement Description} \\
	\midrule
	\endhead
	FR51 & Tech Leads shall be able to configure drift detection thresholds and quality policies. \\
	FR52 & The system shall display real-time documentation health dashboards and freshness metrics. \\
	FR53 & Tech Leads shall be able to manage, version, and preview document templates. \\
	FR54 & Administrators shall manage RBAC permissions, language parsers, and AI settings. \\
	FR55 & Administrators shall monitor background processing pipelines, job statuses, and system health. \\
	\bottomrule
\end{longtable}

\section{Non-Functional Requirements}

\subsection{Performance}
\begin{longtable}{p{2cm} p{11.3cm}}
	\toprule
	\textbf{ID} & \textbf{Requirement Description} \\
	\midrule
	\endhead
	NFR01 & The system shall complete drift analysis for a PR within 2 minutes for repos up to 100k LOC. \\
	NFR02 & The system shall process at least 20 concurrent PR events without performance degradation. \\
	NFR03 & Standard API response times shall remain below 2 seconds under normal workload. \\
	NFR04 & Full-text search queries within the Knowledge Base shall return results within 1 second. \\
	\bottomrule
\end{longtable}

\subsection{Availability \& Scalability}
\begin{longtable}{p{2cm} p{11.3cm}}
	\toprule
	\textbf{ID} & \textbf{Requirement Description} \\
	\midrule
	\endhead
	NFR05 & The system shall maintain an operational uptime availability of at least 99.5\%. \\
	NFR06 & Background worker failures shall not impact user-facing web services. \\
	NFR07 & AI processing workers shall support horizontal auto-scaling based on queue depth. \\
	\bottomrule
\end{longtable}

\subsection{Security \& Reliability}
\begin{longtable}{p{2cm} p{11.3cm}}
	\toprule
	\textbf{ID} & \textbf{Requirement Description} \\
	\midrule
	\endhead
	NFR08 & All API endpoints and communications shall enforce HTTPS/TLS encryption. \\
	NFR09 & Access shall be strictly controlled via Role-Based Access Control (RBAC). \\
	NFR10 & Sensitive credentials and integration keys shall be encrypted at rest. \\
	NFR11 & System audit logs shall be immutable and retention-configurable. \\
	NFR12 & Drift detection algorithms shall yield deterministic evaluations for identical code diffs. \\
	\bottomrule
\end{longtable}

\section{Business Rules}
\begin{longtable}{p{2cm} p{11.3cm}}
	\toprule
	\textbf{ID} & \textbf{Business Rule} \\
	\midrule
	\endhead
	BR01 & All AI-generated or auto-updated documentation must undergo multi-tiered human review before publication. \\
	BR02 & Staff members serve as the first-line technical reviewers. \\
	BR03 & Manager approval is mandatory for documents marked with high sensitivity or critical drift impact. \\
	BR04 & Unapproved documentation drafts shall never be automatically committed to primary repository branches. \\
	BR05 & Critical drift alerts may trigger configurable Pull Request merge blocking policies. \\
	\bottomrule
\end{longtable}

\section{Summary}
This Software Requirements Specification defines the functional capabilities, non-functional constraints, user roles, and operational business rules governing the LivingDocs platform. It serves as the authoritative baseline for architecture design, system implementation, and testing phases.